\documentclass[10pt,a4paper,openany,twoside]{book}
%------------------------------------------------引用的宏包和相应的定义-----------------------------------------
\RequirePackage{silence}
\WarningFilter{fontspec}{Font}
\ActivateWarningFilters

\usepackage{xeCJK}          % 中文支持宏包
\usepackage{newunicodechar}
\usepackage[export]{adjustbox}      

\iffalse
\setCJKmainfont[BoldFont={SimHei},ItalicFont={[simkai.ttf]}]{SimSun}
\setCJKsansfont{SimHei}
\setCJKmonofont{[simfang.ttf]}

\setCJKfamilyfont{zhsong}{SimSun}
\setCJKfamilyfont{zhhei}{SimHei}
\setCJKfamilyfont{zhkai}{[simkai.ttf]}
\setCJKfamilyfont{zhfs}{[simfang.ttf]}
\setCJKfamilyfont{zhli}{LiSu}
%\setCJKfamilyfont{zhyou}{YouYuan}


\newcommand*{\songti}{\CJKfamily{zhsong}}
\newcommand*{\heiti}{\CJKfamily{zhhei}}
\newcommand*{\kaishu}{\CJKfamily{zhkai}}
\newcommand*{\fangsong}{\CJKfamily{zhfs}}
\newcommand*{\lishu}{\CJKfamily{zhli}}
%\newcommand*{\youyuan}{\CJKfamily{zhyou}}
\fi

\usepackage[paperwidth=170mm,paperheight=240mm,left=20mm,right=20mm, top=25mm, bottom=15mm]{geometry}               %定义版面paperwidth=185mm,paperheight=230mm
\usepackage[CJKbookmarks, colorlinks, bookmarksnumbered=true,pdfstartview=FitH,linkcolor=black]{hyperref}   % 选项去掉链接红色方框, 去掉书签可得到想要的印刷效果
%\usepackage{etex}                           %定义更多的计数器,否则警告"no room for a new \count"
%-------------------------------------------------------------------------------------------------------
%\usepackage{graphicx}
%\usepackage{picins}                         %图文混排(文包图),可带方框
%\usepackage{picinpar}                  %图文混排(文包图)  picins在前,picinpar在后,否则有问题
%\usepackage{overpic}                    %加背景图等
%-------------------------------------------------------------------------
\usepackage{fancyhdr}                    % 页眉和页脚的相关定义
\usepackage[perpage]{footmisc}    % 脚注控制
%\usepackage[perpage,symbol]{footmisc}    % 脚注控制
\usepackage{titlesec}                        % 控制标题的宏包
\usepackage{titletoc}                         % 控制目录的宏包
\usepackage{caption2}                       %浮动图形和表格标题样式
\usepackage{setspace}                       % 定制表格和图形的多行标题行距，行间距的宏包
%------------------------------------------------------------------------
\usepackage[usenames,dvipsnames]{color}   % 支持彩色
\usepackage{xcolor}               %单元格加背景
\usepackage{multirow}           %制作表格主要用到\multicolumn，\multirow和\cline,其中，要使用\multirow
%--------------------------------------------------------------------------------------------------------------------------
%\usepackage{everb}                  % 扩展了verbatim，可以加边框、前景色，背景色和行号等，\begin{colorboxed} 这里的内容可以跨页\end{colorboxed}
%\usepackage{listings}          % 可对关键词、注释和字符串等使用不同的字体和颜色，也可以为代码添加边框、背景等风格
\usepackage{fancybox}           % 边框,有阴影,fancybox提供了五种式样\fbox,\shadowbox,\doublebox,\ovalbox,\Ovalbox。
%\usepackage[listings,theorems]{tcolorbox}                % 用于定义类似于beamer的环境框
\usepackage{float}
\usepackage{multicol}           %用于实现在同一页中实现不同的分栏
%\usepackage{flushend}       %分两栏后,\flushend 命令使双栏底部对齐,\raggedend 命令则取消底部对齐,文档最后有效
%\usepackage{balance}           %分两栏后,\balance  命令使双栏底部对齐,\balance   命令则取消底部对齐,任何位置有效

%--------------------------------------------------------------------------------------------
\usepackage{extarrows}          %实现长等号等  This package provides some extensible arrows in maths. mode    (more than \xlongequal \xleftarrow \xrightarrow \xLongleftarrow \xLongrightarrow \xLongleftrightarrow
% wasysym 会在小字号下产生 U/wasy 字形替换警告；本书正文没有使用其特殊符号。
%---------------------------------------------------------------------------------
\usepackage{indentfirst}               % 首行缩进宏包
%\usepackage{color}                        % 支持彩色
%------------------------------------------------------------------------------------
% XeLaTeX 下不要使用旧的 times/helvet 字体包，否则会产生 TU/ptm 与 TU/phv 字形警告。
\IfFontExistsTF{TeX Gyre Termes}
  {\setmainfont{TeX Gyre Termes}
   \setsansfont{TeX Gyre Heros}
   \setmonofont{TeX Gyre Cursor}}
  {}
\usepackage{amssymb}            % AMSLaTeX宏包 用来排出更加漂亮的公式
\usepackage{mathrsfs}               % 不同于\mathcal or \mathfrak 之类的英文花体字体
\usepackage{bm}                         % 处理数学公式中的黑斜体的宏包

\usepackage{tikz,pifont,diagbox}   % 带圆圈的数字，例如\circled{4} 
\newcommand{\circled}[1]{%
  \tikz[baseline=(char.base)]{
    \node[shape=circle,draw,inner sep=1pt,line width=0.5pt] (char) {\small #1};
  }%
}

% 清理日志：旧稿中少量中文标点/文字会落入 Latin 字体或数学模式。
\newunicodechar{，}{\ifmmode\text{\CJKfamily{zhsong}，}\else{\CJKfamily{zhsong}，}\fi}
\newunicodechar{尼}{\ifmmode\text{\CJKfamily{zhsong}尼}\else{\CJKfamily{zhsong}尼}\fi}
\newunicodechar{姆}{\ifmmode\text{\CJKfamily{zhsong}姆}\else{\CJKfamily{zhsong}姆}\fi}
\newunicodechar{≤}{\ensuremath{\leq}}
\hbadness=10000
\vbadness=10000
\hfuzz=80pt
\vfuzz=80pt

\usepackage[amsmath,thmmarks]{ntheorem}         % 定理类环境宏包，其中 amsmath 选项用来兼容 AMS LaTeX 的宏包thmmarks 选项，可以自动恰当地放置定理类环境的结束标记；它还能像图形目录那样生成定理类环境目录
\usepackage{graphicx}
\usepackage{epsfig}
\usepackage{booktabs}
\usepackage{rotating}%将表格或者图片处理成单独一个页面
\usepackage{imakeidx}%可生成索引的标准 LaTeX 宏包
\usepackage{pdfpages}%在LaTeX里插入全页的pdf

%----------------------------------------------------------------------------------------------------------
%\usepackage{flafter}                               % 因为图形可浮动到当前页的顶部， 要防止这种情况，可使用 flafter 宏包
%\usepackage[below]{placeins}
%\usepackage{floatflt}                              % 图文混排用宏包
%\usepackage{rotating}                            % 图形和表格的控制
%\usepackage{type1cm}                           % tex1cm宏包，控制字体的大小
%\usepackage{endfloat}                            %可将浮动对象放置到文件的最后
%----------------------------------------------------------------------------------------------------------------------
\usepackage{fancyref}                                % 支持引用的宏包
%\usepackage{cite}                                    % 支持引用缩写的宏包
%\usepackage{natbib}
\makeatletter
\@ifundefined{CJKnumber}
  {\def\CJKnumber#1{\ifcase#1\or
                    1\or2\or3\or4\or5\or
                    6\or7\or8\or9\or10\or11\or12\or13\or14\or15\or16\or17\or18\or19\or20\fi}}{}
\pdfstringdefDisableCommands{%
  \def\CJKnumber#1{#1}%
  \def\CJKprechaptername{第}%
  \def\thechapter{}%
  \def\CJKthechapter{}%
  \def\CJKchaptername{章}%
  \def\color#1{}%
  \def\hei{}%
  \def\song{}%
}

%\usepackage{amsthm}
 

\renewcommand\contentsname{目\qquad 录}
\renewcommand\listfigurename{插图}
\renewcommand\listtablename{表格}

\@ifundefined{chapter}
  {\renewcommand\refname{参考文献}}
  {\renewcommand\bibname{参考文献}}

\renewcommand\indexname{索引}

\renewcommand\figurename{图}

\renewcommand\tablename{表}

\newcommand\CJKprepartname{第}
\newcommand\CJKpartname{部分}
\newcommand\CJKthepart{\CJKnumber{\@arabic\c@part}}

\@ifundefined{chapter}{}{
  \newcommand\CJKprechaptername{第}
  \newcommand\CJKchaptername{章}
  \newcommand\CJKthechapter{\CJKnumber{\@arabic\c@chapter}}}
\AtBeginDocument{%
  \pdfstringdefDisableCommands{%
    \def\CJKnumber#1{#1}%
    \def\CJKprechaptername{}%
    \def\CJKthechapter{}%
    \def\CJKchaptername{}%
  }%
}
  

\renewcommand\appendixname{附录~\@Alph\c@chapter}

\@ifundefined{mainmatter}
  {\renewcommand\abstractname{摘要}}{}

% \renewcommand\ccname{}                     %   ?
% \renewcommand\enclname{附件}
% \newcommand\prepagename{}                  %   ?
% \newcommand\postpagename{}                 %   ?
% \renewcommand\headtoname{}                 %   ?
% \renewcommand\seename{}                    %   ?

\let\CJK@todaysave=\today
\def\CJK@todaysmall{~\the\year~年~\the\month~月~\the\day~日}
\def\CJK@todaybig{\CJKdigits{\the\year}年\CJKnumber{\the\month}月\CJKnumber{\the\day}日}
\def\CJK@today{\CJK@todaysmall}
\renewcommand\today{\CJK@today}
\newcommand\CJKtoday[1][1]{%
  \ifcase#1\def\CJK@today{\CJK@todaysave}
  \or\def\CJK@today{\CJK@todaysmall}
  \or\def\CJK@today{\CJK@todaybig}
  \fi}

%
% modify the definitions of Standard LaTeX document class
%
\@ifundefined{chapter}{
  \def\@part[#1]#2{%
      \ifnum \c@secnumdepth >\m@ne
        \refstepcounter{part}%
%       \addcontentsline{toc}{part}{\thepart\hspace{1em}#1}%
        \addcontentsline{toc}{part}{\CJKprepartname\expandafter\noexpand\CJKthepart\CJKpartname\hspace{1em}#1}%
      \else
        \addcontentsline{toc}{part}{#1}%
      \fi
      {\parindent \z@ \raggedright
       \interlinepenalty \@M
       \normalfont
       \ifnum \c@secnumdepth >\m@ne
%        \Large\bfseries \partname\nobreakspace\thepart
         \Large\bfseries \CJKprepartname\CJKthepart\CJKpartname
         \par\nobreak
       \fi
       \huge \bfseries #2%
       \markboth{}{}\par}%
      \nobreak
      \vskip 3ex
      \@afterheading}
}{
  \def\@part[#1]#2{%
      \ifnum \c@secnumdepth >-2\relax
        \refstepcounter{part}%
%       \addcontentsline{toc}{part}{\thepart\hspace{1em}#1}%
        \addcontentsline{toc}{part}{\CJKprepartname\expandafter\noexpand\CJKthepart\CJKpartname\hspace{1em}#1}%
      \else
        \addcontentsline{toc}{part}{#1}%
      \fi
      \markboth{}{}%
      {\centering
       \interlinepenalty \@M
       \normalfont
       \ifnum \c@secnumdepth >-2\relax
%        \huge\bfseries \partname\nobreakspace\thepart
         \huge\bfseries \CJKprepartname\CJKthepart\CJKpartname
         \par
         \vskip 20\p@
       \fi
       \Huge \bfseries #2\par}%
      \@endpart}
  \if@twoside
    \def\chaptermark#1{%
      \markboth {\MakeUppercase{%
        \ifnum \c@secnumdepth >\m@ne
          \if@mainmatter
%           \@chapapp\ \thechapter. \ %
            \CJKprechaptername\CJKthechapter\CJKchaptername \ %
          \fi
        \fi
        #1}}{}}%
    \def\sectionmark#1{%
      \markright {\MakeUppercase{%
        \ifnum \c@secnumdepth >\z@
%         \thesection. \ %
          \thesection \ %
        \fi
        #1}}}
  \else
    \def\chaptermark#1{%
      \markright {\MakeUppercase{%
        \ifnum \c@secnumdepth >\m@ne
          \if@mainmatter
%           \@chapapp\ \thechapter. \ %
            \CJKprechaptername\CJKthechapter\CJKchaptername \ %
          \fi
        \fi
        #1}}}
  \fi
  \def\@chapter[#1]#2{\ifnum \c@secnumdepth >\m@ne
                         \if@mainmatter
                           \refstepcounter{chapter}%
%                          \typeout{\@chapapp\space\thechapter.}%
                           \typeout{\CJKprechaptername\CJKthechapter\CJKchaptername}%
                           \addcontentsline{toc}{chapter}%
%                                    {\protect\numberline{\thechapter}#1}%
                                     {\texorpdfstring{\protect\numberline{}\CJKprechaptername%
                                      \expandafter\noexpand\CJKthechapter\CJKchaptername\hspace{0.8em}#1}{#1}}%
                         \else
                           \addcontentsline{toc}{chapter}{#1}%
                         \fi
                      \else
                        \addcontentsline{toc}{chapter}{#1}%
                      \fi
                      \chaptermark{#1}%
                      \addtocontents{lof}{\protect\addvspace{10\p@}}%
                      \addtocontents{lot}{\protect\addvspace{10\p@}}%
                      \if@twocolumn
                        \@topnewpage[\@makechapterhead{#2}]%
                      \else
                        \@makechapterhead{#2}%
                        \@afterheading
                      \fi}
  \def\@makechapterhead#1{%
    \vspace*{50\p@}%
    {\parindent \z@ \raggedright \normalfont
      \ifnum \c@secnumdepth >\m@ne
        \if@mainmatter
%         \huge\bfseries \@chapapp\space \thechapter
          \huge\bfseries \CJKprechaptername\CJKthechapter\CJKchaptername
          \par\nobreak
          \vskip 20\p@
        \fi
      \fi
      \interlinepenalty\@M
      \Huge \bfseries #1\par\nobreak
      \vskip 40\p@
    }}
  \renewcommand*\l@chapter[2]{%
    \ifnum \c@tocdepth >\m@ne
      \addpenalty{-\@highpenalty}%
      \vskip 1.0em \@plus\p@
%     \setlength\@tempdima{1.5em}%
      \setlength\@tempdima{0em}%
      \begingroup
        \parindent \z@ \rightskip \@pnumwidth
        \parfillskip -\@pnumwidth
        \leavevmode \bfseries
        \advance\leftskip\@tempdima
        \hskip -\leftskip
        #1\nobreak\hfil \nobreak\hb@xt@\@pnumwidth{\hss #2}\par
        \penalty\@highpenalty
      \endgroup
    \fi}
 

\let\@appendix\appendix
\renewcommand\appendix{\@appendix%
  \def\CJKprechaptername{\relax}%
  \def\CJKthechapter{\relax}%
  \def\CJKchaptername{\appendixname}}
}  %end of \@ifundefined{chapter}

\def\numberline#1{\ifdim\@tempdima>0pt%
  \settowidth\@tempdimb{#1\space}%
  \ifdim\@tempdima<\@tempdimb%
    \@tempdima=\@tempdimb%
  \fi%
  \hb@xt@\@tempdima{#1\hfil}%
  \fi}
\long\def\@makefntext#1{\leavevmode\hbox{}\kern1em $^{\@thefnmark}$\footnotesize#1}
\fboxsep = 3pt
\fboxrule = .4pt
% Speciali footnotetext pirmam lapui
\def\sfootnotetext[#1]{\begingroup\let\protect\noexpand
      \xdef\@thefnmark{#1}\endgroup
    \@footnotetext}
\makeatother 


\usepackage{natbib}
\bibliographystyle{apalike}

\hypersetup{
  colorlinks=true,
  citecolor=blue,
  linkcolor=blue,
  urlcolor=cyan,
  hypertexnames=false,
  pdftitle={Overleaf Example},
}

\usepackage{amsmath}
\usepackage[most]{tcolorbox}
\newtcolorbox{mybox}[1]{
   colback=white,
    colframe=black,
    colbacktitle=gray!30,
    coltitle=black,
    fonttitle=\bfseries,
    title=#1,
    % --- 关键排版设置 ---
    float=!ht,  % ! 表示强制尝试，“h”当前位置，“t”页顶
}

\newtheorem{exercise}{习题}[chapter]
   %引用宏包所在位置
\makeindex
\usepackage{dcolumn}
\usepackage[UTF8,fontset=none]{ctex}
\setCJKmainfont[
  Script=CustomDefault,
  BoldFont={FandolSong-Bold.otf},
  ItalicFont={FandolKai-Regular.otf}
]{FandolSong-Regular.otf}
\setCJKsansfont[
  Script=CustomDefault,
  BoldFont={FandolHei-Regular.otf}
]{FandolHei-Regular.otf}
\setCJKmonofont[Script=CustomDefault]{FandolFang-Regular.otf}
\setCJKfamilyfont{zhsong}[
  Script=CustomDefault,
  BoldFont={FandolSong-Bold.otf},
  ItalicFont={FandolKai-Regular.otf}
]{FandolSong-Regular.otf}
\setCJKfamilyfont{zhhei}[
  Script=CustomDefault,
  BoldFont={FandolHei-Regular.otf}
]{FandolHei-Regular.otf}
\setCJKfamilyfont{zhkai}[Script=CustomDefault]{FandolKai-Regular.otf}
\setCJKfamilyfont{zhfs}[Script=CustomDefault]{FandolFang-Regular.otf}
\setCJKfamilyfont{zhli}[
  Script=CustomDefault,
  BoldFont={FandolSong-Bold.otf},
  ItalicFont={FandolKai-Regular.otf}
]{FandolSong-Regular.otf}
\usepackage{multicol}
\usepackage[most]{tcolorbox}
\usepackage{graphicx}  
\usepackage{amsmath}
\usepackage{amsfonts}
\usepackage{amssymb}

% 导言区
\graphicspath{{Fig/}{images/}{translation/Fig/}{translation/images/}}% 可选：设置搜索路径(注意大小写)

% --- 为图片文件添加搜索路径 ---
% 告诉LaTeX，图片文件也可以在 ../images/ 目录下找到
% 注意末尾的斜杠 /

% --- 为 \input 和 \include 的 .tex 文件添加搜索路径 ---
% 这需要使用一个内部宏 \input@path
% 你可以把所有可能的 .tex 文件父目录都加进去
\makeatletter
\def\input@path{{../}{../chapters/}}
\makeatother

% --- 为其他文件（如 .sty, .cls, .bib）添加搜索路径 ---
% 这是最通用的方法，但需要编译时设置环境变量

\newcommand{\twodots}{%
  \mathbin{\vcenter{\offinterlineskip
    \hbox{$\cdot$}\vskip1pt\hbox{$\cdot$}}}%
}

\begin{document}

\input{SETUP/FORMAT.tex}    %格式所在位置
  \fontsize{12}{20}\selectfont%
  \baselineskip 20pt %
  \parindent 2em
\pagenumbering{Roman}   %Roman字体书写页码

%-----------------------------------------------------------------------------

% ------------------------------封面--------------------------------------------
\ctitle{\erhao 博弈论基础}

\cauthor{Bernhard von Stengel\hspace{1em}著}


%\ctitle{数学的艺术——第二册\\{\vspace*{0.1em}\emph{在剑桥的茶歇时间}}}

%\cauthor{B\'ela Bollob\'as\hspace{1em}著\\\vspace*{0.2em}张慧铭~~滕广强\hspace{1em}译~~~~\vspace*{0.2em}孙晨旻\hspace{1em}校}
%\publisher{{\hei 世界图书出版公司}\\\vspace*{0.2em} 北京}

%\ctitle{初等概率论——附随机过程和金融数学导论\\\vspace*{0.1em}}
%\cauthor{(第4版)\\\vspace*{0.2em}钟开莱~~Farid AitSahlia\hspace{1em}著\\\vspace*{0.2em}张慧铭\hspace{1em}译}
%\publisher{{\hei 世界图书出版公司}\\\vspace*{0.2em} 北京}

\makecover

% -----------------------------------------------------------------------------------
 %备注：这里文件夹的名字必须严格区分大小写

\centerline{\sihao\hei{\textcolor{black}{XXX}}}\vspace{2cm}






 
\input{PREFACE/PreFACE1.tex} 
%---------------------------------------目录部分---------------------------------
\setcounter{tocdepth}{1}      %重新开始页码


% 1. 手动打印单栏的“目录”标题
% 如果你用的是 ctexbook 或 ctexrep，请用 \chapter*
% 如果你用的是 ctexart，请用 \section*
% \contentsname 是一个会自动变成“目录”或“Contents”的命令
\chapter*{\contentsname} 
%\addcontentsline{toc}{chapter}{\contentsname} % (可选) 让“目录”本身也出现在目录中

% 2. 在两栏环境中，只加载并排版目录条目
\begin{multicols}{2}
    \makeatletter      % 允许使用带@的内部命令
    \@starttoc{toc}    % 这个命令只负责读取.toc文件并生成列表，不包含标题
    \makeatother       % 恢复@的特殊含义
\end{multicols}


\mainmatter    %前言和目录页码结束，正文重新开始设置页码
%-----------------------------------------正文开始---------------------------------
\input{BOOK/chap1}       %第一章
\input{BOOK/chap2}       %第二章
\input{BOOK/chap3}       %第三章
\input{BOOK/chap4}       %第四章
\input{BOOK/chap5}       %第五章
\input{BOOK/chap6}
\input{BOOK/chap7}
\input{BOOK/chap8}       %第八章
\input{BOOK/chap9}       %第八章
\input{BOOK/chap10} 
\input{BOOK/chap11}
\input{BOOK/chap12}
\newpage
\markboth{索引}{索引}
%\clearpage  %清楚前面的偶数空白页
\phantomsection %加这个命令后，目录中的超链接才指向正确的页码
\addcontentsline{toc}{chapter}{\protect\numberline{}{索引}}
\thispagestyle{empty}
\vspace*{2.2cm}

% Page 1
abelian group, 1, 10, 14  阿贝尔群

affine combination, 174  仿射组合

Allais paradox, 121, 133  阿莱悖论 

almost complementary, 251  几乎互补

almost completely labeled, 231, 251, 259  几乎完全标记 

angle, 192, 194  角

antisymmetric, 8  反对称

artificial equilibrium, $230,231,235,250$, 251  {\color{blue} 人工均衡}

associative, 10  结合律

Aumann, Robert, xii, 74, 334   罗伯特·奥曼 

background material, xi  背景材料

\quad converging subsequences, 188  收敛子列

\quad directed graphs and trees, 80  有向图与树图

\quad Euclidean norm and distance, 187  欧几里得范数与距离

\quad expected value, 81  期望值

\quad line segments, 130  线段

\quad matrix multiplication, 141, 143  矩阵乘法

\quad partial orders, 8  偏序 

\quad undirected graphs, 170  无向图

\quad visualizing matrix multiplication, 143  矩阵乘法的可视化 

backward induction, 83, 314, 319  逆向归纳法

\quad more special than SPE, 284, 293 比子博弈完美均衡更特殊

bargaining axioms, 302, 303 讨价还价公理

bargaining solution, 303, 329 讨价还价解

barycenter, 176-178, 190  重心

barycentric coordinates, 176 重心坐标

basic solution, 252  基解

basic variables, 251  基变量

Battle of the Sexes, 60  性别战

behavior strategy, 280, 285  行为策略

best response, 46, 57, 58  最优反应

best-response condition, 135, 145, 148  最优反应条件

best-response diagram, 227  最优反应示意图

bimatrix game, 142  双矩阵博弈

binary system, 1, 5  二进制

bluff, 269 虚张声势

Bolzano-Weierstrass theorem, 188  波尔查诺-魏尔施特拉斯定理

Braess paradox, 43, 44, 49 布雷斯悖论

Brouwer's fixed-point theorem, 151, 153, 189, 195, 203, 258  布劳威尔不动点定理 

castle, 234  城堡

Chicken, 60, 74, 334 胆小鬼博弈

Chomp, 31, 33, 35  {\color{blue} 巧取游戏}

cold game, 30  {\color{blue}  冷静博弈 }

column shift, 239, 240  {\color{blue} 列移位}

commitment game, 93  承诺博弈

common knowledge, 55, 132  共同知识

commutative, 10  交换律

compact set, 149, 151, 171, 190, 212, 248, 302  紧集

complementary, 251  互补性

complementary pivoting, 252, 255, 260  {\color{blue} 互补换基法}

completely labeled, 172,173,181, 227-229, 247, 249, 251  完全标记

concave function, 130  凹函数

concave utility function, 131, 312 凹效用函数

conditional probability, 269, 337  条件概率

Condorcet Paradox, 110, 133  孔多塞悖论

congestion game, 45  拥塞博弈

congestion network, 40, 43, 44  拥塞网络

continuity axiom, 108  连续性公理

Continuous Ultimatum game, 316  连续最后通牒博弈

contractive map, 195  压缩映射

convex combination, 150, 165, 174  凸组合

convex hull, 174, 301  凸包

convex set, 149, 174  凸集



%谭志斌的部分

cooperative game theory, 299, 330  合作博弈论

coordination game, 61, 236 协调博弈

copycat principle, 13 模仿原则

correlated equilibrium, 338, 339 相关均衡

correlated moves, 282, 290 相关行动

correlated strategies, 336 相关策略

correlation device, 337, 338 {\color{blue}相关机制}

cost, 41, 206 成本

Cournot duopoly, 66, 74 古诺双寡头博弈

Cram, 25, 32, 34, 36 {\color{blue} 塞砖棋}

decision theory, 110 决策理论

Dedekind cut, 31 戴德金分割

degenerate game, 162, 164, 233, 238 退化博弈

degree of a map (Brouwer), 203 映射度(布劳威尔)

degree of graph node, 170 图节点的度

descendant, 80, 91 后代节点

deviation, 41, 46 偏离

dictionary, 251 字典

difference trick, 156, 209 差分技巧

digraph, 45 有向图

digraph game, 17, 31, 35 有向图博弈

dimension, 174 维度

dominance solvable, 53, 65, 74, 342 占优可解

dominated 被占优的

\qquad strategy, 342 （被占优）策略

dominated option, 27 被占优选项

Domineering, 24 纵横棋

domino, 24, 32, 34, 36 多米诺骨牌

downward closed set, 313 下闭集

duality of linear programming, 218 线性规划的对偶性

duplicate label, 231 重复标签

edge, 40, 170, 176, 181, 230 边

end node of a path, 170, 232 路径的终止节点

ending condition, 2, 7, 10 终止条件

endpoint of an edge, 170 一条边的端点

entering variable, 252 入基变量

equilibria (plural of equilibrium), 59 均衡（equilibrium 的复数形式）

equilibrium, 53, 58, 73 均衡

\qquad correlated, 338, 339 相关均衡

\qquad in a congestion game, 41, 46, 48 拥塞博弈中的均衡

\qquad mixed, 148, 152 混合均衡

\qquad Nash, 59, 338 纳什均衡

\qquad pure, 53 纯策略均衡

\qquad refinement, 62, 74 均衡精炼

\qquad selection, 62 均衡选择

equilibrium path, 95 均衡路径

equivalence 等价性

\qquad  of impartial games, 11 无偏博弈的等价性

\qquad  of partizan games, 25, 31 有偏博弈的等价性

\qquad  relation, 11 关系的等价性

\qquad  strategic, 239, 259 策略的等价性

\qquad  strong strategic, 241, 259 强策略的等价性

Euclidean distance, 187 欧几里得距离

Euclidean norm, 187 欧几里得范数

expected payoff, 85, 148 期望效用

expected value, 81 期望值

expected-utility function, 103, 106, 112, 117, 124, 208 期望效用函数

exponential size of game tree, 90 博弈树的指数级规模

extreme point, 176, 204 极点

face, 176, 245 面

facet, 176 {\color{blue} 极面}

Farkas’s Lemma, 220 法卡斯定理

first-mover advantage, 96 先行优势

flow, 41 流

\qquad   atomic (non-splittable), 39, 48 原子流(不可拆分)

\qquad   splittable, 40, 48, 50 可拆分流

follower, 96 跟随者

Freudenthal triangulation, 197 弗赖登塔尔三角剖分

full-dimensional, 197 满维

Gambit software, xii Gambit软件 xii

game position, 7 博弈位置

game sum, 10 博弈之和

Game Theory Explorer (GTE), xii

game tree, 3, 79, 270 博弈树

generic game, 84 {\color{blue} 通有博弈}

goalpost method, 159, 208, 226 {\color{blue}  球门柱法 }

graph, 170 图

Grundy value, 16, 31 Grundy值

Hackenbush 哈肯布什游戏

\qquad  Black-White, 30, 37 黑白哈肯布什游戏

\qquad  Gray, 36 灰色哈肯布什游戏

halfline, 192 {\color{blue} 射线}

halfspace, 193, 245 半空间

heap, see Nim heap 堆，尼姆堆 

history of moves, 3, 81 行动的历史

%朱峰的部分

Homo Ludens, (Huizinga) 《游戏的人》(赫伊津哈 著), 330

hot game, 30  {\color{blue} 热情博弈 }

hotel with seaview, 188, 204 {\color{blue}  海景酒店 }

hyperplane, 193 超平面 

identity matrix, 344 单位矩阵 

impartial game, 2 无偏博弈  

imperfect recall, 275 不完美记忆  

incentive constraints, 338 激励约束 

independence axiom, 104, 108, 121 独立性公理

index of an equilibrium, 259 均衡的指数

indifference, 106, 112, 139 无差异 

information set, 265, 266 信息集 

Inspection game, 137 监管博弈 

insurance 保险 

\begin{itemize}
\item[]accidental breakdown, 109 意外故障 

\item[]car rental, 109, 110 汽车租赁 
\end{itemize}

intransitive preferences, 110, 114, 133 非传递偏好 

invariant under utility scaling, 303 效用缩放不变性 

irrelevant alternatives, 303 无关选项

iterated elimination of dominated, 64, 65, 69, 71, 222 重复剔除劣势策略 


iterated-offers bargaining game, 317-330 重复讨价还价博弈


Kayles, 17, 23, 31  凯尔斯博弈 


Kuhn's theorem, 284 库恩定理 

label, 172, 227 标签 

Last Year in Marienbad, 4 《去年在马里昂巴德》

leader, 96 领导  

leadership game, 96 领导博弈 

leaf, 80 叶节点 

learning, 349 学习 

leaving variable, 252 离基变量 

Left player, 24 左玩家 

Lemke's algorithm, 260 莱姆基算法 

Lemke-Howson 莱姆基 - 豪森
\begin{itemize}
\item[] algorithm, 232  算法 
\item[] graph, 233, 235, 259 图 
\item[] path, 231, 232, 244 路径
\end{itemize} 

lexico-minimum ratio test, 258 字典序最小比率检验 

lexico-positive, 257 字典正序 

lexicographic order, 105, 114, 123, 124 字典序 

LH, see Lemke-Howson 莱姆克 - 豪森

line segment, 130 线段 

linear complementarity problem, 260 线性互补问题 

linear programming, 218, 221, 255, 260, 296  线性规划 

linear programming duality, 206, 218 线性规划对偶 

losing position, 3, 25 失败位置

lottery, 106, 111 彩票 
\begin{itemize}
\item[] simple, 104, 117 简单  
\end{itemize} 

lower envelope, 209 下包络 

L-position, 25 {\color{blue}左胜位置} 

majority vote, 110 {\color{blue}多数人规则} 

Marienbad, 4 马里昂巴德 

Markov chain, 344 马尔可夫链 

Maschler, Michael, 334 迈克尔·马希勒 

Matching Pennies, 72 硬币匹配（博弈） 

matrix game, 207 矩阵博弈 

matrix multiplication, 141, 143 矩阵乘法 

max-min strategy, 138, 210, 211 极大极小策略 

maximizer, 207 {\color{blue}最大值点} 

mex rule, 15, 16, 31 mex 规则 

min-max strategy, 211, 212 极小极大策略 

minimal element, 8 极小元 

minimax theorem, 213, 214, 216, 344, 347 极小极大定理 

minimizer, 207 {\color{blue}最小值点} 

minimum ratio test, 255 最小比值检验 

misère Nim, 4, 32 {\color{blue}反向尼姆游戏} 

misère play, 2, 14 {\color{blue}反向玩法} 

missing label, 231 缺失标签 

mixed equilibrium, 148, 152, 228 混合均衡 

mixed extension, 139, 148, 212 混合扩展 

mixed strategy, 143 混合策略 

mixed-strategy profile, 148 混合策略组合 

mixed-strategy set, 144 混合策略集 

negative game, 11, 13, 25 负博弈 

neighbor, 170 邻居 

Nim, 2, 32 尼姆游戏

Nim heap, 2, 15 尼姆堆 

Nim sum, 5, 20 尼姆和 

Nim value, 16, 27 尼姆值 

nimber, 27 尼姆数 

node, 170, 230 节点 

nonbasic variables, 252 非基变量

nondegenerate, 254 非退化的 

normal form, 74 标准型 


Page 4  %--（王思杰）

normal play, 2 正常规则(与反常规则相对应)


normal vector, 193, 245 法向量


normative, 112, 114, 120, 133 规范的


Northcott's game, 31, 34 诺思科特博弈


N-position, 25 先手胜位置


number (partizan game), 27, 28 数(有偏博弈中的)


optimal strategy, 215 最优策略


option, 3, 7, 10, 26 选项

partial order , 8 偏序

total order, 8 全序

origin of utility scale, 104, 107, 11 效用尺度原点


orthogonal, 191 正交的


outcome class, 3, 25 结果类


outcome set for decisions, 111 决策的结果集


parallel edges, 40, 41, 49 平行边


Pareto-optimal, 303 帕累托最优


parity argument, 170, 184, 203, 232, 234 奇偶性论证



partial strategy profile, 339 部分策略组合


partizan game,2,24,37,38 有偏博弈


path, 45, 80, 170 路径


payoff equivalent, 62 收益等价


perfect information, 2 完美信息


perfect recall, 275-279, 285 完美记忆


permutation, 198 排列

Pigou network, 40, 42, 44, 49 庇古网络


pivoting, 252 换基操作


play of a game, 77 博弈的进行


Poker, 268, 297 扑克


Poker Nim, 15, 16, 17 扑克尼姆游戏


polyhedron, 245, 341 多面体


polytope, 175, 245 有界多面体


potential function, 39, 46, 48, 49 势函数


potential game, 50, 74 势博弈


PPAD, 259 (Polynomial Parity Argument(s) on Directed Graphs)有向图上的多项式时间奇偶性论证类

P-position, 25 后手胜位置


preference 偏好


relation, 111 关系

strict, 106, 112 严格的

Price of Anarchy, 49 无政府代价


prime number decomposition theorem, 9 质数分解定理


Prisoner's Dilemma, 55, 73, 300 囚徒困境


product graph, 231 {\color{blue} 乘积图}


projection, 171, 191 投影


Quality game, 62 产品质量博弈


quantity competition, 66 产量竞争


Queen-move game, 22, 34 皇后移动博弈


random variable, 81 随机变量


rationality, 66, 132 理性

sequential, 91 顺序的


ray, 192 射线


ray termination, 255 射线终止

realization equivalence, 284 实现等价性


recursive, 4, 10 递归的


reduced strategic form, 87, 275 精简策略型


reduced strategy, 78, 87, 274 精简策略


reflexive, 8, 11 自反的


regret, 109, 122, 349 后悔


rental car insurance, 109 租车保险


Right player, 24 右玩家


Rock-Paper-Scissors, 72, 74 石头剪刀布


Rook-move game, 17 战车移动游戏


R-position, 25 右胜位置


scalar product, 142, 144 标量积


scale-equivalent, 115, 116, 118 尺度等价的


Schelling, Thomas, 74 托马斯·谢林


seaview hotel, 188, 204 {\color{blue}海景酒店}


security strategy, 212 安全策略


Seoul (Korea), 44, 49 首尔(韩国)


separating hyperplane, 193, 220 分离超平面


sequence form, 295 序列形式


set inclusion, 8 集合包含


sg-value, 16 sg值


simple lottery, 104, 117 简单彩票


simplex, 175, 339 单纯形


simplex algorithm, 255, 260 单纯形算法


simplicial subdivision, 176 单纯剖分


slack variables, 251 松弛变量


Social Choice, 110, 133 社会选择


social optimum, 41, 42, 49 社会最优

SPE, see subgame-perfect equilibrium 子博弈完美均衡(SPE)


Sperner condition,180,184 斯佩纳条件


Sperner labeling, 180 斯佩纳标号


Sprague-Grundy value, 16, 31 斯普拉格 - 格鲁迪值


Stackelberg game, 96, 98 斯塔克尔伯格博弈


Page 5

Stag Hunt, 73 猎鹿博弈  

staggered payoffs, 56, 74 {\color{blue}交错收益  }

stationary distribution, 344 平稳分布  

stationary strategy, 325 平稳策略  

strategic equivalence, 239 策略等价  

strategic form, 74, 273 策略形式  

strict order, 8, 32 严格序  

strict preference, 106, 112 严格偏好  

strictly concave, 131 严格凹  

strictly dominated strategy, 62 严格被占优策略  

strong strategic equivalence, 241 强策略等价  

sub-facet, 176, 178, 182 {\color{blue}子极面}  

sub-segment, 173 子线段  

sub-simplex, 176 子单纯形  

subgame, 91, 291 子博弈  

subgame-perfect equilibrium, 91, 292 子博弈完美均衡  

subjective probability, 111 主观概率  

subtree, 91, 291 子树  

symmetric 对称的  

\quad bimatrix game, 142 对称双矩阵博弈  

\quad equilibrium, 61, 73 对称均衡  

\quad game, 55, 60, 66, 72 对称博弈  

symmetry axiom, 303 对称性公理  

temperature scale, 104, 107, 116 温度刻度  

terminal node, 80 终端节点  

threat point, 302 威胁点  

tight inequality, 245 紧不等式  

top-down induction, 8 {\color{blue} 自上而下归纳  }

total order, 8 全序  

total relation, 112 全关系  

tradeoff, 114 权衡  

transformation 变换  

\quad positive-affine, 115 正仿射变换  

\quad projective, 249 射影变换  

\quad strictly increasing, 115 严格递增变换  

transitive, 8, 11 传递的  

transposed matrix, 141 转置矩阵  

triangulation, 176 三角剖分  

Trust Dilemma, 61 信任困境  

Ultimatum game, 314 最后通牒博弈  

unit cube, 196 单位立方体  

unit of utility scale, 104, 107, 116 效用刻度单位  

unit simplex, 151, 169–171, 175, 190 单位单纯形  

unit vector, 150, 169, 171 单位向量  

upper envelope, 158, 225, 241, 246 上包络  

utility function, 103–133 效用函数  

\quad cardinal, 104, 107, 115 基数效用函数  

\quad ordinal, 106, 115 序数效用函数  

utility scale, 104, 116, 137, 240 效用尺度  

vertex, 176, 204, 245 顶点  

walk, 80 行走  

Wardrop equilibrium, 50 {\color{blue}用户均衡 } 

weakly dominated strategy, 62, 63, 64, 75, 219 弱被占优策略  

\quad elimination may lose equilibria, 63, 220 剔除弱劣势策略可能导致均衡遗漏  

Wilson, Robert, 259 罗伯特·威尔逊

winning move, 3, 4 必胜步  

winning position, 3, 25 必胜局势  

zero-sum game, 207 零和博弈  





\footnotesize{
\printindex
}       %第三章


\iffalse
%\include{preface/intro}
%\include{preface/intro2}

\centerline{\sihao\hei{\textcolor{black}{XXX}}}\vspace{2cm}






                                      %译者序
\include{preface/preface1.TEX}                               %4版前言
%\include{preface/preface2}                                     %金融前言
%\include{preface/preface3}                               %3版前言
%\include{preface/preface4}                               %2版前言
%\include{preface/preface5}                               %1版前言
\include{preface/c_abstract}                             %中文摘要
\include{preface/e_abstract}                             %英文摘要

\include{preface/preface1}                               %序言，译者序
\renewcommand{\thefigure}{\arabic{figure}}
\include{book/chap1}       %第一章
\include{book/chap2}       %第二章
\include{book/chap3}       %第三章
%\include{book/chap4}       %第四章
%\include{book/appendix1}
%\include{book/chap5}       %第五章
%\include{book/chap6}       %第六章
%\include{book/chap7}       %第七章
%\include{book/appendix2}
%\include{book/chap8}       %第八章
%\include{book/appendix3}
%\include{book/chap9}       %第九章
%\include{book/appendix4}
%\include{book/chap10}       %第十章
%\include{book/chap11}   %第十一章
%\include{book/chap12}   %第十二章
%-----------------------------------------正文结束--------------------------------
%\include{reference/references}                              %参考文献部分
%\include{book/answer}                                          %习题答案
%\include{book/appendix5}
%\include{book/index}                                          %索引
\fi


\end{document}
